\documentclass[conference,letterpaper,10pt]{IEEEtran}
\usepackage{amsmath}
\usepackage{amssymb}
\usepackage{graphicx}
\usepackage{booktabs}
\usepackage{siunitx}
\usepackage{hyperref}
\usepackage{array}

\title{Dynamic Localization Across Space and Time\\
A Frequency-Based Framework for Spatial Position, Temporal Re-Localization, and Informational Retrieval}
\author{
  Christopher Woodyard\\
  Vers3Dynamics, R.A.I.N. Lab\\
  Washington, D.C., USA
}

\begin{document}
\maketitle

\begin{abstract}
This paper is a spatial localization and temporal re-localization under a single resonance model in which
\textbf{location is a property of the object}, selected by resonance between matter and a background scalar field.
Rather than treating position and time as external labels on a spacetime manifold, we encode the spatiotemporal
position of any physical system in the dominant resonance frequency \(\omega_{\rm loc}\) of its coupled
matter-scalar-field state. We define a frequency-parameterized localization operator \(\mathcal{L}\) acting on the
extended Hilbert space \(\mathcal{H}_{\rm sys} \otimes \mathcal{H}_\phi\), derive the driven field-theoretic coupling,
and recover standard quantum mechanics and classical trajectories in the weak-coupling limit. Macroscopic material
re-localization (spatial or temporal) remains energetically prohibitive (\(E_{\rm req} > 10^{45}\) J) while
informational re-localization respects Landauer's bound and is physically feasible. The framework supplies concrete,
falsifiable predictions for atomic-clock interferometry, matter-wave interferometry, and gravitational-wave detection.
\end{abstract}

\begin{IEEEkeywords}
Dynamic Location Theory, scalar fields, frequency-location operator, temporal resonance, emergent spacetime, quantum
supercomputing, Landauer's principle, falsifiable quantum field theory, gravitational-wave detection
\end{IEEEkeywords}

\section{Introduction}
In standard physics, position and time serve as external labels on a spacetime manifold. While this framework has been
 extraordinarily successful, it does not provide a \emph{dynamical explanation} for why an object occupies one
position rather than another, nor why temporal evolution proceeds in one direction versus another. Dynamic Location
Theory (DLT) proposes that localization itself is an emergent \emph{property of the object}, selected by resonance
between matter and a background scalar field. Spatial position and temporal coordinate become two manifestations of
the same frequency-tuning problem.
This distinction matters because it separates two categories of re-localization that standard physics conflates:
material transport (moving an object through space or time) and informational access (retrieving a state that was
localized at a different spatiotemporal coordinate). DLT recovers all standard predictions in appropriate limits,
renders material transport energetically forbidden, and remains falsifiable through concrete experimental predictions.

\section{Unified Theoretical Framework}
\subsection{Localization Operator}
We define the unified resonance operator \(\mathcal{L}\) acting on the tensor-product Hilbert space \(\mathcal{H}_{\rm
sys} \otimes \mathcal{H}_\phi\):
\begin{equation}
\mathcal{L}\,|\Psi\rangle = \omega_{\rm loc}\,|\Psi\rangle, \qquad
|\Psi\rangle = |\psi\rangle_{\rm matter} \otimes |\phi\rangle_{\rm scalar},
\label{eq:loc-op}
\end{equation}
where \(\omega_{\rm loc}\) is the dominant resonance signature of the coupled matter-scalar system. The scalar field
\(|\phi\rangle_{\rm scalar}\) describes the background field state, whose vacuum expectation value
\(\langle\phi\rangle_0\) parameterizes the classical spacetime background.
\textbf{Central Thesis (Formal).} \emph{Location is a property of the object:} the observable spatiotemporal
coordinates \((x,t)\) of a physical system are eigenvalues of the localization operator \(\mathcal{L}\) acting on its
coupled state, not pre-existing labels on a manifold.

\subsection{Unified Hamiltonian}
The total non-relativistic Hamiltonian governing the coupled system is
\begin{equation}
H_{\rm total} = H_{\rm matter} + H_{\rm scalar} + \gamma\hbar\bigl(\omega_w - \omega_{\rm loc}(\mathbf{x},t)\bigr),
\label{eq:hamiltonian}
\end{equation}
where \(\gamma\) is a dimensionless coupling constant and \(\omega_w\) is the driving frequency of an external field.
Forward (adiabatic) evolution corresponds to the tracking condition \(\omega_w \approx \omega_{\rm loc}\);
re-localization to a target resonance requires non-adiabatic driving.

\subsection{Three Modes of Re-Localization}
Table~\ref{tab:modes} classifies the distinct physical meanings and energy costs of re-localization under DLT.
\begin{table*}[t]
\centering
\small
\caption{Classification of re-localization modes by domain, physical meaning, and energy scale.}
\label{tab:modes}
\begin{tabular}{@{}p{1.5cm} p{2.0cm} p{6.5cm} p{4.0cm}@{}}
\toprule
\textbf{Mode} & \textbf{Domain} & \textbf{Physical Meaning} & \textbf{Energy Scale} \\
\midrule
Type I & Parallel / branch & Phase shift between resonant states & Low to moderate \\
Type II & Material & Displacement of mass-energy & Prohibitive (\(>10^{45}\) J) \\
Type III & Informational & Memory and data retrieval & Near Landauer's limit \\
\bottomrule
\end{tabular}
\end{table*}
The energetic separation between Type II and Type III is central to the framework: moving macroscopic matter (Type II)
 is distinguished not categorically but quantitatively from retrieving stored information (Type III).

\section{Spatial Localization as a Dynamical Variable}
The frequency-location operator in position representation is
\begin{equation}
\mathcal{L} = \int d^3x\;\omega(\mathbf{x})\,|\phi(\mathbf{x})\rangle\langle\phi(\mathbf{x})| \otimes
\hat{P}_{\mathbf{x}},
\label{eq:spatial-op}
\end{equation}
where \(\hat{P}_{\mathbf{x}} = |\mathbf{x}\rangle\langle\mathbf{x}|\) projects onto position eigenstates. In the uniform
 vacuum ground state \(|\phi_0\rangle\), the scalar field factorizes and we recover the standard position operator as a
marginal case.
Weak-coupling corrections to standard position probabilities are
\begin{equation}
\epsilon(\mathbf{x},t) = \xi^2\,\frac{|\phi(\mathbf{x},t)-\phi_0|^2}{\phi_0^2}\,
f\!\left(\frac{\omega_{\rm loc}(\mathbf{x})}{\omega_w}\right) \sim 10^{-26},
\label{eq:epsilon}
\end{equation}
for laboratory values \(\xi \sim 10^{-10}\), yielding corrections at the \(10^{-26}\) level---well beyond current
experimental resolution but theoretically well-defined.

\section{Temporal Re-Localization as the Same Mechanism}
Temporal eigenstates are selected by the \emph{same} resonance condition. Forward time evolution corresponds to the
adiabatic tracking \(\omega_w \approx \omega_{\rm loc}(t)\). Re-localization to \(t' \neq t\) is non-adiabatic resonance
tuning to a different temporal eigenfrequency.
Concretely, a state \(|\psi(t)\rangle\) localized at time \(t\) can be driven to a state \(|\psi(t')\rangle\) by applying a
driving field at frequency \(\omega_w = \omega_{\rm loc}(t')\) for a finite duration \(\Delta t\). The energy cost scales
with the overlap \(\langle\psi(t')|\psi(t)\rangle\) and the driving amplitude, not with the temporal separation \(|t'-t|\)
 itself---this is why informational re-localization (Type III) is energetically distinct from material temporal
transport (Type II).
The distinction between ``time travel'' and memory access is therefore purely energetic, not categorical: both are
resonance tuning problems differing only in the complexity of the target state.

\section{Field-Theoretic Description}
The effective action incorporating frequency-dependent non-minimal coupling between the scalar field and gravity is
\begin{equation}
S = \int d^4x\sqrt{-g}\left[\frac{R}{16\pi G} - \frac12(\partial\phi)^2 - V(\phi) - \xi(\omega)\phi^2 R +
\mathcal{L}_{\rm SM}\right],
\label{eq:action}
\end{equation}
where \(R\) is the Ricci scalar, \(G\) is Newton's gravitational constant, and \(\mathcal{L}_{\rm SM}\) is the Standard
Model Lagrangian. The frequency-dependent coupling \(\xi(\omega)\) encodes the resonance response:
\begin{equation}
\xi(\omega) = \xi_0\!\left(1 + \alpha\, e^{-(\omega-\omega^*)^2/\Delta\omega^2}\right), \qquad \xi_0 \sim 10^{-10}.
\label{eq:xi}
\end{equation}
Here \(\omega^*\) is the resonance center, \(\Delta\omega\) is the resonance width, and \(\alpha\) is a dimensionless
amplitude. The Gaussian form arises naturally from the Breit-Wigner resonance profile of the coupled matter-scalar
system.
The driven Klein-Gordon equation in flat spacetime (\(g^{\mu\nu} = \eta^{\mu\nu}\)) follows from
varying~\eqref{eq:action}:
\begin{equation}
\ddot{\phi} - \nabla^2\phi + m^2\phi \approx J_0\,e^{-i\omega_w t}\,\delta^3(\mathbf{x}-\mathbf{x}_0).
\label{eq:klein-gordon}
\end{equation}
The scalar field exerts an effective force on matter via the gradient of the coupling energy density:
\begin{equation}
\mathbf{F}_\phi = -\nabla\!\bigl(\xi(\omega)\,\phi^2\,\rho_{\rm matter}\bigr) = -\nabla U_\phi(\mathbf{x},t).
\label{eq:force}
\end{equation}
This force derives from the potential \(U_\phi = \xi(\omega)\phi^2\rho_{\rm matter}\), confirming energy conservation in
 the coupled system.

\section{Classical and Quantum Limits}
\subsection{Classical Limit (Ehrenfest Theorem)}
Applying the Ehrenfest theorem to the localization operator yields the classical trajectory in the limit \(\hbar \to
0\):
\begin{equation}
\lim_{\hbar\to0}\langle\mathcal{L}(t)\rangle = \omega\bigl(\mathbf{x}_{\rm cl}(t)\bigr).
\label{eq:classical}
\end{equation}
The classical equation of motion follows from the Hellmann-Feynman force:
\begin{equation}
m\ddot{\mathbf{x}}_{\rm cl} = \mathbf{F}_{\rm classical} + \mathbf{F}_\phi.
\label{eq:eom}
\end{equation}
\subsection{Weak-Coupling Limit (Newtonian Mechanics)}
In the limit of vanishing scalar-matter coupling (\(\xi \to 0\)), the scalar force vanishes and the trajectory reduces
to the Newtonian solution:
\begin{equation}
\lim_{\xi\to0}\mathbf{x}_{\rm cl}(t) = \mathbf{x}_{\rm Newton}(t).
\label{eq:newton}
\end{equation}
Thus DLT contains Newtonian mechanics as a limiting case, as required for any viable physical theory.

\section{Energy Constraints}
Material displacement (Type II re-localization, spatial or temporal) obeys the scaling
\begin{equation}
E_{\rm req} \approx mc^2\,\xi^2\!\left(\frac{\Delta}{\lambda_0}\right)^{\!2},
\label{eq:energy}
\end{equation}
where \(\Delta\) is the spatial separation or normalized temporal separation \(\Delta t/\tau_P\), \(\lambda_0 = \hbar/(mc)\)
 is the Compton wavelength, and \(\tau_P = \sqrt{\hbar G/c^5}\) is the Planck time. For a macroscopic object (\(m \sim
1\;\text{kg}\), \(\Delta \sim 1\;\text{m}\)), this yields \(E_{\rm req} \gg 10^{45}\) J, rendering material re-localization
 energetically prohibitive.
Informational re-localization (Type III), by contrast, satisfies Landauer's principle:
\begin{equation}
E_{\rm info} \ge k_B T\ln 2 \approx 10^{-21}\;\text{J/bit}\quad (T = 300\;\text{K}).
\label{eq:landauer}
\end{equation}
The 24-order-of-magnitude gap between~\eqref{eq:energy} and~\eqref{eq:landauer} for macroscopic systems is the
framework's central empirical prediction.

\section{Experimental Considerations and Falsifiability}
\subsection{Predictions}
\textbf{Prediction 1 (Atomic-Clock Interferometry).} Differential interferometry between two spatially separated
atomic clocks measures a relative phase shift
\begin{equation}
\Delta\phi \approx \gamma\int_{t_0}^{t_1}\!\bigl(\omega_1(t) - \omega_2(t)\bigr)\,dt,
\label{eq:phase}
\end{equation}
arising from the coupling \(\gamma\hbar(\omega_w - \omega_{\rm loc})\). With next-generation optical clocks achieving
fractional uncertainty \(\Delta\nu/\nu \sim 10^{-19}\), this yields phase shifts of order \(10^{-19}\) rad in a
\SI{1}{\second} integration time.
\textbf{Prediction 2 (Matter-Wave Interferometry).} A matter-wave interferometer exhibits anomalous phase shifts when
the scalar-field gradient exceeds
\begin{equation}
|\nabla\phi|^2 > 10^{16}\;\text{(SI units)}.
\label{eq:gradient-threshold}
\end{equation}
as derived from equating the scalar force \(F_\phi\) with the interferometer's intrinsic phase sensitivity.
\textbf{Prediction 3 (Gravitational-Wave Interferometry).} Gravitational-wave detectors (LIGO, Virgo, KAGRA, LISA)
operating near the resonance center \(\omega^*\) will exhibit a narrow-band resonant excess strain or
frequency-dependent deviation from the General Relativistic antenna pattern, with amplitude scaling as \(\alpha\xi_0
h\), where \(h\) is the gravitational-wave strain amplitude.

\subsection{Falsification Criteria}
The framework is falsified if any of the following are observed:
\begin{enumerate}
\item Atomic-clock interferometry achieves \(10^{-20}\) rad precision with no phase noise correlating with \(\omega_{\rm
loc}\) differences between the clocks.
\item Matter-wave interferometers detect no anomalous phase shift at gradient sensitivities down to \(10^{15}\) SI
units.
\item Gravitational-wave detectors show no frequency-dependent excess strain or antenna-pattern deviation at the level
 \(\gamma\xi_0 \sim 10^{-10}\) near the predicted resonance frequency.
\item High-fidelity quantum simulations of the \(\mathcal{H}_{\rm sys}\otimes\mathcal{H}_\phi\) coupled system exhibit
instabilities that violate energy or momentum conservation at the \(10^{-6}\) level.
\end{enumerate}

\section{Type I Re-Localization: Mathematical Derivation}
Type I re-localization corresponds to a coherent phase shift between two (or more) resonant eigenstates of the
localization operator \(\mathcal{L}\) without any net displacement of mass-energy. Start from the localization operator
eigenvalue equation
\begin{equation}
\mathcal{L}\,|\Psi_i\rangle = \omega_i\,|\Psi_i\rangle, \qquad i=1,2,
\label{eq:type1-eigen}
\end{equation}
where \(|\Psi_i\rangle = |\psi_i\rangle_{\rm matter} \otimes |\phi_i\rangle_{\rm scalar}\).
The interaction piece of the total Hamiltonian is
\begin{equation}
H_{\rm int} = \gamma\hbar\,(\omega_w\,\mathbb{I} - \mathcal{L}).
\label{eq:type1-hint}
\end{equation}
Projecting onto the two-dimensional subspace spanned by \(\{|\Psi_1\rangle,|\Psi_2\rangle\}\), the effective Hamiltonian
matrix is
\begin{equation}
H_{\rm eff} = \gamma\hbar\begin{pmatrix}
\omega_w - \omega_1 & 0 \\
0 & \omega_w - \omega_2
\end{pmatrix}.
\label{eq:type1-heff}
\end{equation}
The unitary time-evolution operator is
\begin{equation}
U(t) = \exp\!\Bigl(-i\gamma\int_0^t(\omega_w(t')-\mathcal{L})\,dt'\Bigr).
\label{eq:type1-evo}
\end{equation}
Acting on a coherent superposition \(|\Psi(0)\rangle = \alpha|\Psi_1\rangle + \beta|\Psi_2\rangle\) yields
\begin{equation}
|\Psi(t)\rangle = \alpha\,e^{-i\gamma\int_0^t(\omega_w(t')-\omega_1)\,dt'}|\Psi_1\rangle
+ \beta\,e^{-i\gamma\int_0^t(\omega_w(t')-\omega_2)\,dt'}|\Psi_2\rangle.
\label{eq:type1-superpos}
\end{equation}
The relative phase between the two branches is
\begin{equation}
\Delta\phi(t) = \gamma\int_0^t\bigl(\omega_1(t') - \omega_2(t')\bigr)\,dt'.
\label{eq:type1-deltaphi}
\end{equation}
This is precisely the phase shift measured in differential atomic-clock interferometry (Prediction 1). The energy cost
 of Type I re-localization is
\begin{equation}
E_{\rm I} \approx \gamma\hbar\,\Delta\omega \ll mc^2\xi^2,
\label{eq:type1-energy}
\end{equation}
many orders of magnitude below the Type II threshold and comparable to or below Landauer's limit.

\section{Application to Gravitational-Wave Detection}
Gravitational-wave detectors (LIGO, Virgo, KAGRA, LISA) are ultra-sensitive Michelson interferometers that measure
differential phase shifts from spacetime strain \(h_{\mu\nu}\). Under DLT, the non-minimal coupling \(-\xi(\omega)\phi^2
R\) sources scalar excitations \(\delta\phi\) from propagating GW curvature. The linearized driven Klein-Gordon equation
in curved spacetime yields a resonant scalar response when the GW frequency lies near \(\omega^*\).
The induced differential resonance-frequency shift between orthogonal arms is
\begin{equation}
\Delta\omega_{\rm loc}(t) = \gamma\,\xi(\omega)\,h_+(t)\,\phi_0^2,
\label{eq:gw-deltaomega}
\end{equation}
where \(h_+\) is the plus-polarized GW strain and \(\phi_0 = \langle\phi\rangle_0\) is the vacuum scalar field amplitude.
The resulting DLT phase shift in the interferometer output is
\begin{equation}
\Delta\phi_{\rm DLT} = C\,\xi(\omega)\,h_+(t)\,\tau,
\label{eq:gw-phase}
\end{equation}
where \(\tau = 2L/c\) is the round-trip light travel time for arm length \(L\), and the dimensionless constant
\begin{equation}
C = \gamma\,\phi_0^2\,(2L/c)
\label{eq:gw-constant}
\end{equation}
absorbs the coupling strength, scalar field amplitude, and interferometer geometry. This adds coherently to the
standard General Relativistic strain signal, producing a narrow-band resonant excess strain peaked at \(\omega^*\) with
amplitude scaling as \(\alpha\xi_0 h\).
The entire GW-induced signal is a pure Type I relative phase shift: no test-mass displacement (Type II) occurs, and
data analysis constitutes informational re-localization (Type III) governed by Landauer's bound.

\section{Conclusion}
Dynamic Location Theory establishes that \textbf{location is a property of the object} by identifying the dominant
resonance frequency \(\omega_{\rm loc}\) of the coupled matter-scalar system as the fundamental localization variable.
This single principle unifies spatial position, temporal re-localization, and memory under one mathematical framework,
 recovers standard quantum and classical predictions in appropriate limits, and makes concrete falsifiable predictions
 accessible to near-term experiments in atomic-clock, matter-wave, and gravitational-wave interferometry. Material
transport remains energetically forbidden; informational access is not.

\section*{Acknowledgments}
This work was conducted independently at Vers3Dynamics and the R.A.I.N. Lab, Washington, D.C.

\begin{thebibliography}{10}
\bibitem{woodyard2026location}
C. Woodyard, ``Location Is a Dynamic Variable: A Scalar-Field Framework for Emergent Localization,'' Vers3Dynamics,
Jan. 2026, doi: 10.5281/zenodo.18263032.
\bibitem{woodyard2026temporal}
C. Woodyard, ``Temporal Re-Localization via Scalar Resonance,'' Vers3Dynamics, R.A.I.N. Lab, Jan. 2026, doi:
10.5281/zenodo.18281960.
\bibitem{epr}
A. Einstein, B. Podolsky, and N. Rosen, ``Can quantum-mechanical description of physical reality be considered
complete?'' \emph{Phys. Rev.}, vol. 47, no. 10, pp. 777--780, 1935.
\bibitem{bell}
J. S. Bell, ``On the Einstein-Podolsky-Rosen paradox,'' \emph{Physics}, vol. 1, no. 3, pp. 195--200, 1964.
\bibitem{morris}
M. S. Morris, K. S. Thorne, and U. Yurtsever, ``Wormholes, time machines, and the weak energy condition,'' \emph{Phys.
Rev. Lett.}, vol. 61, no. 13, pp. 1446--1449, 1988.
\bibitem{zurek}
W. H. Zurek, ``Decoherence, einselection, and the quantum origins of the classical,'' \emph{Rev. Mod. Phys.}, vol. 75,
no. 3, pp. 715--775, 2003.
\bibitem{landauer}
R. Landauer, ``Irreversibility and heat generation in the computing process,'' \emph{IBM J. Res. Dev.}, vol. 5, no. 3,
pp. 183--191, 1961.
\bibitem{birrell}
N. D. Birrell and P. C. W. Davies, \emph{Quantum Fields in Curved Space}. Cambridge, U.K.: Cambridge Univ. Press,
1982.
\bibitem{damour}
T. Damour and G. Esposito-Farèse, ``Tensor-multi-scalar theories of gravitation,'' \emph{Class. Quantum Grav.}, vol.
9, no. 9, pp. 2093--2176, 1992.
\bibitem{feynman}
R. P. Feynman, ``Simulating physics with computers,'' \emph{Int. J. Theor. Phys.}, vol. 21, no. 6/7, pp. 467--488,
1982.
\end{thebibliography}

\end{document}
